%%% FIELD stmt-published-occupancy
Under \(\mathrm{ass:iid\text{-}sampling}\), \(\mathrm{ass:overlap}\), \(\mathrm{ass:mean\text{-}envelope}\), \(\mathrm{ass:second\text{-}moment\text{-}envelope}\), and \(\mathrm{ass:maximal\text{-}effect\text{-}radius}\), the clipped matched-cell occupancy estimator \(\widetilde\tau\) computed from the first \(m_0\) observations satisfies, for \(n\geq 8\),
\[
 E_{Q_P^{(n)}}\!\left[(\widetilde\tau-\tau(P))^2\right]
 \leq C_{\epsilon}M^2\left(\sigma^2+\frac1{m_0}+\frac d{m_0^2}\right)
\]
uniformly over \(P\in\mathcal P_{n,d,\epsilon,M,\sigma}\).  The statement is the rescaled real-outcome, bounded-conditional-second-moment form of the published matched-cell occupancy bound.

%%% FIELD stmt-accepted-parent
Throughout the declared range, there are total observed-sample estimators and constants depending only on \(\epsilon\) for which
\[
 c_{\epsilon}M^2\left\{b_{n,d}+\sigma^2\min(1,u_{n,d})\right\}
 \leq \mathsf R_{n,d,\epsilon,M,\sigma}
 \leq C_{\epsilon}M^2\left[\frac1n+\min\left\{1,u_{n,d},\sigma^2+\frac d{n^2}\right\}\right].
\]

%%% FIELD stmt-empirical
Let \(\widehat\tau^{\mathrm{em}}\) be exactly the estimator in \(\mathrm{def:empirical\text{-}mass\text{-}estimator}\), whose pilot is exactly \(\mathrm{def:occupancy\text{-}pilot}\). There is \(C_{\epsilon}<\infty\) such that, uniformly over \(n\ge3\), \(2\le d\le n\), \(0\le\sigma\le2\), and \(P\in\mathcal P_{n,d,\epsilon,M,\sigma}\),
\[
 E_{Q_P^{(n)}}[(\widehat\tau^{\mathrm{em}}-\tau(P))^2]
 \le C_{\epsilon}M^2\left(\frac1n+\sigma^2\frac{d^2}{n^2}\right).
\]

%%% FIELD stmt-bracket
There are constants \(0<c_{\epsilon}\le C_{\epsilon}<\infty\) such that throughout the declared parameter range,
\[
 c_{\epsilon}M^2\{b_{n,d}+\sigma^2\min(1,u_{n,d})\}
 \le\mathsf R_{n,d,\epsilon,M,\sigma}
 \le C_{\epsilon}M^2q_{n,d,\sigma}^{\mathrm{em}}.
\]
This is an infimum-risk upper bound: for \(d\le n\) it combines \(\mathrm{lem:empirical\text{-}mass\text{-}upper}\) with the inherited bounds, and for \(d>n\) it equals the accepted parent's upper envelope. No claim about an algebraically unspecified selector is part of the theorem. Relative to the accepted parent's upper rate \(n^{-1}+\min\{1,u_{n,d},\sigma^2+d/n^2\}\), it is a strict tightening whenever \(d\le n\) and \(d/n^2+\sigma^2d^2/n^2<\min\{u_{n,d},\sigma^2+d/n^2\}\).

%%% FIELD proof-empirical
Write \(m=m_1\), \(t=\widetilde\tau\), \(e=t-\tau(P)\), and
\[
 \delta_k=\tau_k-\tau(P),\qquad
 U=\frac1m\sum_{k=1}^dN_k^{(1)}(1-J_k^{(1)}),\qquad
 S=\frac1m\sum_{k=1}^dN_k^{(1)}\delta_k.
\]
All three quantities are defined even when cells or blocks are empty by \(\mathrm{def:occupancy\text{-}pilot}\) and \(\mathrm{def:empirical\text{-}mass\text{-}estimator}\).  The mean envelope gives \(|\tau_k|\leq M\) on occupied cells and hence \(|\tau(P)|\leq M\).  The maximal-radius assumption gives \(|\delta_k|\leq\sigma M\), while the definition of \(\tau(P)\) gives
\[
 \sum_{k=1}^dp_k\delta_k
 =\sum_{k=1}^dp_k\tau_k-\tau(P)\sum_{k=1}^dp_k=0.
\]

We first prove the needed unmatched-mass estimate.  Put
\[
 s_{ak}=P(X=k,A=a),\qquad
 V_{ak}=N_{ak}^{(1)}\mathbf 1\{N_{1-a,k}^{(1)}=0\}.
\]
Then \(mU=\sum_{k,a}V_{ak}\) exactly: an unmatched nonempty cell contributes its entire occupancy in its unique observed arm, and an empty or matched cell contributes zero.  Fix \((a,k)\), set \(\alpha=s_{ak}\) and \(\beta=s_{1-a,k}\), and use the multinomial factorial-moment identities supplied by iid sampling.  Since \(N^2=N(N-1)+N\),
\[
 E[V_{ak}^2]
 =m\alpha(1-\beta)^{m-1}
  +m(m-1)\alpha^2(1-\beta)^{m-2}.
\]
For distinct cells \(k\ne\ell\), if \(\gamma=s_{c\ell}\) and \(\phi=s_{1-c,\ell}\), the same identity gives
\[
 E[V_{ak}V_{c\ell}]
 =m(m-1)\alpha\gamma(1-\beta-\phi)^{m-2},
\]
whereas \(V_{0k}V_{1k}=0\).  Overlap implies
\(
 \alpha\leq\epsilon^{-1}\beta
\)
and
\(
 \gamma\leq\epsilon^{-1}\phi
\);
these inequalities also hold when the relevant cell mass is zero.  For \(m\geq4\), \((1-x)^r\leq e^{-rx}\) and the finite bounds on \(\sup_{x\geq0}xe^{-x}\) and \(\sup_{x\geq0}x^2e^{-x}\) therefore yield
\[
 E[V_{ak}^2]\leq C_{\epsilon},
 \qquad
 E[V_{ak}V_{c\ell}]\leq C_{\epsilon}\quad(k\ne\ell).
\]
There are \(2d\) diagonal terms and fewer than \(4d^2\) ordered cross terms.  Consequently,
\[
 E[U^2]\leq C_{\epsilon}\frac{d^2}{m^2}.
\]
The deterministic inequality \(0\leq U\leq1\) supplies the uniform refinement
\[
 E[U^2]\leq C_{\epsilon}\min\left\{1,\frac{d^2}{m^2}\right\}. \tag{1}
\]

Next isolate the correction-block outcome noise
\[
 W=\frac1m\sum_{k=1}^dN_k^{(1)}J_k^{(1)}
 \left\{(\overline Y_{1k}^{(1)}-\mu_{1k})
 -(\overline Y_{0k}^{(1)}-\mu_{0k})\right\}.
\]
Let \(\mathscr G\) be generated by the pilot and all correction-block \((X,A)\) values.  Iid sampling makes the outcome residuals conditionally independent across observations, and their conditional means vanish.  Thus \(E[W\mid\mathscr G]=0\), and the second-moment envelope gives
\[
 E[W^2\mid\mathscr G]
 \leq \frac{M^2}{m^2}\sum_{k=1}^d
 (N_k^{(1)})^2J_k^{(1)}
 \left(\frac1{N_{1k}^{(1)}}+\frac1{N_{0k}^{(1)}}\right). \tag{2}
\]
The displayed reciprocals occur only when \(J_k^{(1)}=1\), so no division by zero is present.  Conditional on \(N_k^{(1)}=r\), the treated count is \(B\sim\operatorname{Bin}(r,\pi_k)\).  For \(B\sim\operatorname{Bin}(r,q)\) with \(q>0\),
\[
 E\left[\frac{\mathbf1\{B>0\}}B\right]
 \leq2E\left[\frac1{B+1}\right]
 =\frac{2\{1-(1-q)^{r+1}\}}{(r+1)q}
 \leq\frac2{(r+1)q}. \tag{3}
\]
Apply (3) to the two arms, discard the other-arm indicator, and use \(\pi_k,1-\pi_k\geq\epsilon\).  The conditional expectation of the \(k\)-th summand in (2) is at most
\[
 \frac{4r^2}{(r+1)\epsilon}\leq\frac{4r}{\epsilon}.
\]
Since \(\sum_kE[N_k^{(1)}]=m\), (2) yields
\[
 E[W^2]\leq\frac{4M^2}{\epsilon m}. \tag{4}
\]

On a matched correction cell write
\(
 \widehat\Delta_k^{(1)}=\overline Y_{1k}^{(1)}-\overline Y_{0k}^{(1)}
\), and put
\[
 T_{\rm raw}=t+\frac1m\sum_{k=1}^dN_k^{(1)}J_k^{(1)}
 (\widehat\Delta_k^{(1)}-t).
\]
The value assigned to \(\widehat\Delta_k^{(1)}\) off the event \(J_k^{(1)}=1\) is immaterial because its coefficient is zero.  Before the final clipping, direct algebra in the estimator definition gives the exact identity
\[
 T_{\rm raw}-\tau(P)
 =W+S+\frac1m\sum_{k=1}^dN_k^{(1)}(1-J_k^{(1)})(e-\delta_k). \tag{5}
\]
Indeed, expand \(T_{\rm raw}=t+m^{-1}\sum_kN_k^{(1)}J_k^{(1)}(\widehat\Delta_k^{(1)}-t)\), insert \(\widehat\Delta_k^{(1)}=\tau_k\) plus the residual defining \(W\) on matched cells, and use \(\sum_kN_k^{(1)}=m\).  Because \(m^{-1}\sum_{i=m_0+1}^n\delta_{X_i}=S\), iid sampling and the centering identity above imply
\[
 E[S]=0,\qquad E[S^2]\leq\frac{\sigma^2M^2}{m}. \tag{6}
\]
Moreover, if the last term of (5) is denoted by \(H\), then
\[
 |H|\leq U(|e|+\sigma M). \tag{7}
\]
The pilot and correction blocks are independent by iid sampling.  Hence (1) and (7) imply
\[
 E[H^2]
 \leq2\{E[e^2]+\sigma^2M^2\}E[U^2]. \tag{8}
\]
The cited matched-cell occupancy bound \(\mathrm{lem:published\text{-}occupancy\text{-}risk}\), applied to the first \(m_0\) records, gives
\[
 E[e^2]\leq C_{\epsilon}M^2
 \left(\sigma^2+\frac1{m_0}+\frac d{m_0^2}\right).
\]
For \(n\geq8\), \(m_0,m_1\geq n/2-1\); since \(d\leq n\), this reduces to
\[
 E[e^2]\leq C_{\epsilon}M^2(\sigma^2+n^{-1}),
 \qquad
 E[U^2]\leq C_{\epsilon}d^2/n^2. \tag{9}
\]
The conditional mean-zero property of \(W\) makes its cross term with the other terms in (5) vanish.  Applying \((x+y)^2\leq2x^2+2y^2\) to \(S+H\), and then using (4), (6), (8), and (9), gives
\[
 \begin{aligned}
 E[(T_{\rm raw}-\tau(P))^2]
 &\leq E[W^2]+2E[S^2]+2E[H^2]\\
 &\leq C_{\epsilon}M^2\left\{
 \frac1n+\frac{\sigma^2}{n}
 +(\sigma^2+n^{-1})\frac{d^2}{n^2}\right\}\\
 &\leq C_{\epsilon}M^2\left(\frac1n+\sigma^2\frac{d^2}{n^2}\right),
 \end{aligned}
\]
where the last line uses \(d\leq n\) and \(0\leq\sigma\leq2\).  Since \(\tau(P)\in[-M,M]\), projection onto \([-M,M]\) cannot increase squared error, so the same bound holds for \(\widehat\tau^{\mathrm{em}}\).  Finally, when \(3\leq n<8\), the definition returns zero and the mean envelope gives \((0-\tau(P))^2\leq M^2\leq8M^2/n\).  Enlarging \(C_{\epsilon}\) proves the claim uniformly over every stated sample size, including null cells and entirely unmatched blocks.

%%% FIELD proof-bracket
The accepted same-class result \(\mathrm{lem:accepted\text{-}parent\text{-}bracket}\) gives, on exactly \(\mathcal P_{n,d,\epsilon,M,\sigma}\),
\[
 c_{\epsilon}M^2\{b_{n,d}+\sigma^2\min(1,u_{n,d})\}
 \leq\mathsf R_{n,d,\epsilon,M,\sigma}
\]
and supplies a particular total data-driven estimator with risk at most
\[
 C_{\epsilon}M^2\left[\frac1n+
 \min\left\{1,u_{n,d},\sigma^2+\frac d{n^2}\right\}\right]. \tag{10}
\]
This proves the asserted lower inequality without changing the model class.

For the upper inequality, denote the total estimator supplied with (10) by \(T^{\rm par}\).  If \(d>n\), then \(v_{n,d}=1\), and the right side of (10) is exactly
\[
 C_{\epsilon}M^2\left[\frac1n+
 \min\left\{1,u_{n,d},\frac d{n^2}+\sigma^2v_{n,d}\right\}\right]
 =C_{\epsilon}M^2q_{n,d,\sigma}^{\mathrm{em}}. \tag{11}
\]
If \(d\leq n\), then \(v_{n,d}=d^2/n^2\), and \(\mathrm{lem:empirical\text{-}mass\text{-}upper}\) gives the second, total, data-driven candidate
\[
 E[(\widehat\tau^{\mathrm{em}}-\tau(P))^2]
 \leq C_{\epsilon}M^2\left(\frac1n+\sigma^2v_{n,d}\right). \tag{12}
\]
Define one estimator from the sample by using \(\widehat\tau^{\mathrm{em}}\) when \(d\leq n\) and
\[
 \sigma^2v_{n,d}
 \leq\min\left\{1,u_{n,d},\sigma^2+\frac d{n^2}\right\},
\]
and using \(T^{\rm par}\) otherwise.  This rule depends only on the known inputs \((n,d,\sigma)\); both branches are total on every sample.  Combining (10) and (12), and enlarging the comparison constant to the larger branch constant, gives
\[
 \begin{aligned}
 \mathsf R_{n,d,\epsilon,M,\sigma}
 &\leq C_{\epsilon}M^2\left[\frac1n+
 \min\left\{1,u_{n,d},\sigma^2+\frac d{n^2},\sigma^2v_{n,d}\right\}\right]\\
 &\leq C_{\epsilon}M^2\left[\frac1n+
 \min\left\{1,u_{n,d},\frac d{n^2}+\sigma^2v_{n,d}\right\}\right]\\
 &=C_{\epsilon}M^2q_{n,d,\sigma}^{\mathrm{em}}.
 \end{aligned}
\]
This proves the upper inequality and exhibits the deterministic selector that witnesses it.

It remains to verify the strict-improvement clause.  When \(d\leq n\),
\(
 u_{n,d}\leq1/\log^2(en)<1
\).
If
\[
 \frac d{n^2}+\sigma^2\frac{d^2}{n^2}
 <\min\left\{u_{n,d},\sigma^2+\frac d{n^2}\right\},
\]
then the third entry in the definition of \(q_{n,d,\sigma}^{\mathrm{em}}\) is strictly smaller than every active entry of the parent's minimum.  Adding the common term \(n^{-1}\) preserves strict inequality, proving the final assertion.

%%% FIELD proof-collision
Assume \(2\leq d\leq n\) and \(\sigma^2d^2\leq n\).  The empirical-mass lemma gives
\[
 \mathsf R_{n,d,\epsilon,M,\sigma}
 \leq C_{\epsilon}M^2\left(\frac1n+\sigma^2\frac{d^2}{n^2}\right)
 \leq \frac{2C_{\epsilon}M^2}{n}.
\]
The lower inequality in \(\mathrm{thm:refined\text{-}bracket}\), together with \(b_{n,d}=n^{-1}+dn^{-2}\geq n^{-1}\), gives
\[
 \mathsf R_{n,d,\epsilon,M,\sigma}\geq c_{\epsilon}M^2/n.
\]
These two inequalities prove the uniform comparison \(\mathsf R_{n,d,\epsilon,M,\sigma}\asymp_{\epsilon}M^2/n\).

For \(d=\lfloor n^{3/4}\rfloor\) and \(\sigma^2=u_{n,d}\), one has \(2\leq d\leq n\) for \(n\geq3\), and
\[
 \sigma^2d^2
 =\frac{d^4}{n^2\log^2(en)}
 \leq\frac{n}{\log^2(en)}\leq n.
\]
Thus this sequence lies in the just-proved region, and its minimax MSE is of order \(M^2/n\), with constants depending only on \(\epsilon\).

%%% FIELD proof-global
Fix \(n\geq3\), \(h\in\{-1,1\}\), and \(z\in\mathcal Z_n\).  Elementary differentiation shows that \(x\mapsto\log(ex)/\sqrt{x}\) decreases for \(x\geq3\), so
\[
 0<\sigma_n^\star\leq\frac{\log(3e)}{\sqrt3}<\frac43.
\]
Also \(\ell_n^\star>1\).  Therefore every effect in \(\mathrm{def:benchmark\text{-}marked\text{-}family}\) satisfies
\[
 \frac{|\tau_k|}{M}
 \leq\frac{\Delta_n^\star}{M}+\frac{\sigma_n^\star}{2}
 =\frac{\sigma_n^\star}{4\ell_n^\star}+\frac{\sigma_n^\star}{2}
 <1. \tag{13}
\]
Thus all displayed Bernoulli probabilities are in \([0,1]\).  Their conditional means are
\[
 E[Y(a)\mid X=k]
 =\frac M2\left[2\left\{\frac12+\frac{(2a-1)\tau_k}{2M}\right\}-1\right]
 =\frac{(2a-1)\tau_k}{2}=\mu_{ak},
\]
which obey the mean envelope by (13), and their conditional central variances are at most \(M^2/4\), hence obey the second-moment envelope.  The construction draws \(A\) independently of both potential outcomes conditional on \(X\) and sets \(Y=Y(A)\); consequently conditional exchangeability and consistency hold.  Its propensity is \(1/2\), which obeys overlap because \(0<\epsilon<1/2\).

Exact mass normalization follows from \(p_k=1/n\) for \(1\leq k\leq n\).  Since \(z\in\mathcal Z_n\) has \(\sum_kz_k=0\),
\[
 \tau(P_{h,z}^\star)
 =\frac1n\sum_{k=1}^n\left(h\Delta_n^\star+\frac{M\sigma_n^\star}{2}z_k\right)
 =h\Delta_n^\star, \tag{14}
\]
and, for every occupied cell,
\[
 |\tau_k-\tau(P_{h,z}^\star)|
 =\frac{M\sigma_n^\star}{2}|z_k|
 \leq\frac{M\sigma_n^\star}{2}
 \leq M\sigma_n^\star. \tag{15}
\]
Equations (13)--(15) and the preceding checks prove \(P_{h,z}^\star\in\mathcal P_{n,n,\epsilon,M,\sigma_n^\star}\).  They also prove that both mixture hypotheses have the identical deterministic design \((p_k,\pi_k)=(1/n,1/2)\), and that their ATE supports are the singletons \(\{h\Delta_n^\star\}\).

For the distance calculation set \(Z_i=(2A_i-1)Y_i\).  Conditional on any fixed \((h,z)\), the \(Z_i\) are iid, lie in \([-M/2,M/2]\), and satisfy
\[
 E[Z_i\mid h,z]
 =\sum_{k=1}^n\frac1n\sum_{a=0}^1\frac12(2a-1)\mu_{ak}
 =\frac12\tau(P_{h,z}^\star)
 =\frac{h\Delta_n^\star}{2}. \tag{16}
\]
For completeness, let a random variable \(V\) lie in \([a,b]\), put \(p=(EV-a)/(b-a)\), and set \(t=\lambda(b-a)\).  Convexity of the exponential gives
\[
 E e^{\lambda(V-EV)}
 \leq e^{-pt}\{1-p+pe^t\}.
\]
If \(g(t)=\log[e^{-pt}\{1-p+pe^t\}]\), then \(g(0)=g'(0)=0\), while \(g''(t)\) is the variance of a Bernoulli variable under an exponential tilt and hence is at most \(1/4\).  Integrating this inequality twice, for either sign of \(t\), gives \(g(t)\leq t^2/8\).  With \(b-a=M\),
\[
 E\exp\{\lambda(V-EV)\}\leq\exp(\lambda^2M^2/8). \tag{17}
\]
Exponential Markov inequality applied to (16), optimized at \(\lambda=2\Delta_n^\star/M^2\), therefore yields uniformly in \(z\)
\[
 Q_{P_{1,z}^\star}^{(n)}(\overline Z_n\leq0)
 \leq\exp\left\{-\frac{n(\Delta_n^\star)^2}{2M^2}\right\},
 \qquad
 Q_{P_{-1,z}^\star}^{(n)}(\overline Z_n\geq0)
 \leq\exp\left\{-\frac{n(\Delta_n^\star)^2}{2M^2}\right\}. \tag{18}
\]
Integrating (18) over the two uniform priors preserves both bounds.  With \(B=\{\overline Z_n>0\}\), the variational definition of total variation then gives
\[
 \begin{aligned}
 \operatorname{TV}(\mathbb Q_{-1}^\star,\mathbb Q_1^\star)
 &\geq \mathbb Q_1^\star(B)-\mathbb Q_{-1}^\star(B)\\
 &\geq1-2\exp\left\{-\frac{n(\Delta_n^\star)^2}{2M^2}\right\}\\
 &=1-2\exp\left\{-\frac{\log^2(en)}{32\log^2(e+\log^2(en))}\right\},
 \end{aligned}
\]
where the last equality uses the definitions of \(\sigma_n^\star\), \(\ell_n^\star\), and \(\Delta_n^\star\).  The exponent tends to infinity because \(\log(en)/\log(e+\log^2(en))\to\infty\).  Hence the lower bound on total variation tends to one and in particular exceeds \(1/3\) for all sufficiently large \(n\).  The full-mixture indistinguishability condition fails despite exact normalization, common design, and realization-wise radius control; equation (16) identifies the hypothesis-wide outcome shift as the cause.

%%% FIELD oeq-partial
The proved results give only an outer bracket on the benchmark separation, not the requested optimization.  Let \(\rho_n^{\mathrm{outer}}\) denote the supremum separation over prior pairs satisfying only exact normalization, supportwise membership in the benchmark class, identical design marginals, target-support interval separation, and mixture total variation at most \(1/3\), without imposing a particular centered-score representation.  At \(d=n\) and \(\sigma=\sigma_n^\star\), \(\mathrm{thm:refined\text{-}bracket}\) and the third branch of \(q_{n,d,\sigma}^{\mathrm{em}}\) imply
\[
 \mathsf R_{n,n,\epsilon,M,\sigma_n^\star}
 \leq C_{\epsilon}M^2(2n^{-1}+(\sigma_n^\star)^2)
 \leq3C_{\epsilon}M^2(\sigma_n^\star)^2. \tag{19}
\]
If two feasible priors have target-support intervals separated by \(M\rho_n\), classify an estimate according to the nearer interval.  A classification error forces absolute estimation error at least \(M\rho_n/2\).  The sum of the two testing errors is at least one minus the total variation of the mixture laws.  Therefore the average squared error is at least
\[
 \frac{M^2\rho_n^2}{8}(1-\operatorname{TV})
 \geq\frac{M^2\rho_n^2}{12}. \tag{20}
\]
Combining (19)--(20) gives the necessary outer bound
\[
 \rho_n\leq C_{\epsilon}\sigma_n^\star
 =C_{\epsilon}\frac{\log(en)}{\sqrt n}. \tag{21}
\]

There is also a matching parametric-order construction for the outer constraints.  Fix an absolute \(0<c\leq1/3\).  For \(h\in\{-1,1\}\), take \(p_k=1/n\), \(\pi_k=1/2\), and \(\tau_k=hcM/\sqrt n\) for every cell.  Put \(\mu_{ak}=(2a-1)\tau_k/2\), and use the same two-point potential-outcome construction as in \(\mathrm{def:benchmark\text{-}marked\text{-}family}\).  Every support law is exactly homogeneous, so its radius is zero; the mean, variance, overlap, consistency, and exchangeability checks are the same as those in \(\mathrm{lem:global\text{-}shift\text{-}prior\text{-}fails}\).  The two ATEs are \(\pm cM/\sqrt n\), with gap \(2cM/\sqrt n\).

To check distance, put \(t=c/(2\sqrt n)\).  Conditional on either arm, the two hypotheses have Bernoulli mark probabilities \(1/2+t\) and \(1/2-t\), whose Hellinger affinity is \(\sqrt{1-4t^2}\).  The design is identical, so the one-observation affinity has the same value and the \(n\)-observation affinity is
\[
 \mathcal A_n=(1-4t^2)^{n/2}. \tag{22}
\]
For any two dominated laws, Cauchy--Schwarz applied to
\(
 |p-q|=|\sqrt p-\sqrt q|(\sqrt p+\sqrt q)
\)
gives \(\operatorname{TV}\leq\sqrt{1-\mathcal A_n^2}\).  Bernoulli's inequality and (22) now yield
\[
 \operatorname{TV}
 \leq\sqrt{1-(1-4t^2)^n}
 \leq2\sqrt n\,t=c\leq\frac13. \tag{23}
\]
Thus this singleton-prior pair certifies order \(n^{-1/2}\) under the outer support, normalization, common-design, target-separation, and distance constraints.  It does not implement the conditionally centered marked-score mechanism.  Consequently the currently proved information is only
\[
 2c n^{-1/2}\leq\rho_n^{\mathrm{outer}}
 \leq C_{\epsilon}\log(en)n^{-1/2},
\]
while \(\mathrm{lem:global\text{-}shift\text{-}prior\text{-}fails}\) proves that the displayed nonparametric global-shift candidate is infeasible.  No realization-wise centered marked coupling attaining a separation larger than order \(n^{-1/2}\), and no universal parametric obstruction for that centered program, follows from the frozen primitives.
